% TikZ-Zeichnung: Arduino Nano 33 BLE Sense
% Eigene schematische Darstellung, angelehnt an das Pinout im Arduino-Datenblatt.
% Kompilierbar als eigenständige Datei.


\begin{tikzpicture}[
	font=\sffamily,
	board/.style={
		rounded corners=2.5mm,
		fill=teal!65!black,
		draw=teal!25!black,
		line width=0.9pt
	},
	part/.style={
		rounded corners=1mm,
		fill=black!80,
		draw=black,
		line width=0.4pt
	},
	module/.style={
		rounded corners=1mm,
		fill=gray!25,
		draw=gray!70!black,
		line width=0.4pt
	},
	sensor/.style={
		rounded corners=0.8mm,
		fill=orange!25,
		draw=orange!70!black,
		line width=0.35pt
	},
	pin/.style={
		circle,
		fill=black!85,
		draw=white,
		line width=0.25pt,
		minimum size=2.1mm,
		inner sep=0pt
	},
	pinlabel/.style={
		font=\scriptsize\sffamily,
		text=black
	},
	powerlabel/.style={
		font=\scriptsize\sffamily,
		text=red!75!black
	},
	analoglabel/.style={
		font=\scriptsize\sffamily,
		text=blue!75!black
	},
	digitallabel/.style={
		font=\scriptsize\sffamily,
		text=violet!80!black
	},
	commlabel/.style={
		font=\scriptsize\sffamily,
		text=green!45!black
	},
	whitetext/.style={
		font=\scriptsize\sffamily,
		text=white
	}
]

% --------------------------------------------------------------------
% Grundform
% --------------------------------------------------------------------
\draw[board, blur shadow={shadow blur steps=5, shadow xshift=0.8mm, shadow yshift=-0.8mm}]
	(-2.35,-5.8) rectangle (2.35,5.8);

% USB-Buchse
\draw[rounded corners=0.8mm, fill=gray!35, draw=gray!70!black, line width=0.6pt]
	(-0.85,5.8) rectangle (0.85,6.35);
\draw[rounded corners=0.4mm, fill=black!75, draw=black]
	(-0.55,5.92) rectangle (0.55,6.22);
\node[font=\tiny\sffamily, text=black] at (0,6.52) {Micro-USB};

% Pinleisten
\foreach \y in {5.05,4.35,3.65,2.95,2.25,1.55,0.85,0.15,-0.55,-1.25,-1.95,-2.65,-3.35,-4.05,-4.75} {
	\node[pin] at (-2.05,\y) {};
	\node[pin] at (2.05,\y) {};
}

% Bauteile auf dem Board
\draw[module] (-1.25,2.55) rectangle (1.25,3.75);
\node[font=\scriptsize\sffamily, align=center] at (0,3.15) {NINA-B306\\nRF52840};

\draw[part] (-0.95,0.85) rectangle (0.95,1.55);
\node[whitetext, align=center] at (0,1.20) {Cortex-M4F\\64 MHz};

\draw[sensor] (-1.35,-0.05) rectangle (-0.25,0.55);
\node[font=\tiny\sffamily, align=center] at (-0.80,0.25) {IMU};

\draw[sensor] (0.25,-0.05) rectangle (1.35,0.55);
\node[font=\tiny\sffamily, align=center] at (0.80,0.25) {APDS};

\draw[sensor] (-1.35,-1.00) rectangle (-0.25,-0.40);
\node[font=\tiny\sffamily, align=center] at (-0.80,-0.70) {LPS};

\draw[sensor] (0.25,-1.00) rectangle (1.35,-0.40);
\node[font=\tiny\sffamily, align=center] at (0.80,-0.70) {HTS};

\draw[part] (-0.60,-2.00) rectangle (0.60,-1.35);
\node[whitetext, align=center] at (0,-1.68) {MIC};

\draw[part] (-1.05,-3.20) rectangle (1.05,-2.55);
\node[whitetext, align=center] at (0,-2.88) {3.3 V Regler};

\draw[fill=yellow!80!orange, draw=yellow!60!black] (-0.38,-4.15) rectangle (0.38,-3.55);
\node[font=\tiny\sffamily] at (0,-3.85) {LED};

% Board-Beschriftung
\node[font=\bfseries\small\sffamily, text=white, align=center] at (0,4.55) {Arduino Nano\\33 BLE Sense};
\node[font=\tiny\sffamily, text=white] at (0,-5.35) {3.3 V I/O \quad | \quad nicht 5-V-tolerant};

% --------------------------------------------------------------------
% Pinbeschriftungen links
% --------------------------------------------------------------------
\node[digitallabel, anchor=east] at (-2.28,5.05) {D13 / SCK / LED};
\node[powerlabel,  anchor=east] at (-2.28,4.35) {+3V3};
\node[analoglabel, anchor=east] at (-2.28,3.65) {AREF};
\node[analoglabel, anchor=east] at (-2.28,2.95) {A0};
\node[analoglabel, anchor=east] at (-2.28,2.25) {A1};
\node[analoglabel, anchor=east] at (-2.28,1.55) {A2};
\node[analoglabel, anchor=east] at (-2.28,0.85) {A3};
\node[commlabel,  anchor=east] at (-2.28,0.15) {A4 / SDA};
\node[commlabel,  anchor=east] at (-2.28,-0.55) {A5 / SCL};
\node[analoglabel, anchor=east] at (-2.28,-1.25) {A6};
\node[analoglabel, anchor=east] at (-2.28,-1.95) {A7};
\node[powerlabel,  anchor=east] at (-2.28,-2.65) {VUSB / 5V*};
\node[pinlabel,    anchor=east] at (-2.28,-3.35) {RST};
\node[powerlabel,  anchor=east] at (-2.28,-4.05) {GND};
\node[powerlabel,  anchor=east] at (-2.28,-4.75) {VIN};

% --------------------------------------------------------------------
% Pinbeschriftungen rechts
% --------------------------------------------------------------------
\node[commlabel,    anchor=west] at (2.28,5.05) {D12 / MISO};
\node[commlabel,    anchor=west] at (2.28,4.35) {D11 / MOSI};
\node[digitallabel, anchor=west] at (2.28,3.65) {D10 / PWM};
\node[digitallabel, anchor=west] at (2.28,2.95) {D9 / PWM};
\node[digitallabel, anchor=west] at (2.28,2.25) {D8};
\node[digitallabel, anchor=west] at (2.28,1.55) {D7};
\node[digitallabel, anchor=west] at (2.28,0.85) {D6 / PWM};
\node[digitallabel, anchor=west] at (2.28,0.15) {D5 / PWM};
\node[digitallabel, anchor=west] at (2.28,-0.55) {D4};
\node[digitallabel, anchor=west] at (2.28,-1.25) {D3 / PWM};
\node[digitallabel, anchor=west] at (2.28,-1.95) {D2};
\node[powerlabel,   anchor=west] at (2.28,-2.65) {GND};
\node[pinlabel,     anchor=west] at (2.28,-3.35) {RST};
\node[commlabel,    anchor=west] at (2.28,-4.05) {RX / D0};
\node[commlabel,    anchor=west] at (2.28,-4.75) {TX / D1};

% --------------------------------------------------------------------
% Legende
% --------------------------------------------------------------------
\begin{scope}[shift={(0,-6.65)}]
	\node[font=\scriptsize\sffamily] at (-3.35,0) {\textcolor{red!75!black}{Power}};
	\node[font=\scriptsize\sffamily] at (-2.10,0) {\textcolor{blue!75!black}{Analog}};
	\node[font=\scriptsize\sffamily] at (-0.75,0) {\textcolor{violet!80!black}{Digital/PWM}};
	\node[font=\scriptsize\sffamily] at (1.00,0) {\textcolor{green!45!black}{Kommunikation}};
	\node[font=\tiny\sffamily, align=center] at (0,-0.45)
		{* VUSB/5V ist abhängig von der VUSB-Lötbrücke und der USB-Versorgung. I/O-Pins bleiben 3,3-V-Pins.};
\end{scope}

\end{tikzpicture}

\end{document}
